% !PRINTED 278
\weberpage{278}
%<<CONT
is, we set
\weq{14}{U = uvv'v'', \quad U_1 = u_1v_1v_1'v_1'', \quad U_2 = u_2v_2v_2'v_2'',}
\weq{15}{
\begin{array}{rl}
2A & = U + (-1)^\mu(U_1+U_2), \\
B & = UU_1 + UU_2 + (-1)^\mu U_1U_2, \\
& = \dfrac{4}{U_1} + \dfrac{4}{U_2} + (-1)^\mu\dfrac{4}{U}.
\end{array}
}

From (12) and § 73, (4) we then obtain the following associated
permutations:
\weq{16}{
\begin{array}{lll}
\omega, & A, & B, \\
-\dfrac{1}{\omega}, & \varrho^{n'+1}A, & \varrho^{-(n'+1)}B, \\
\omega+1, & \sigma^{n'+1}e^{\frac{2\mu\pi i}{3}}A, & \sigma^{-(n'+1)}e^{-\frac{2\mu\pi i}{3}}B.
\end{array}
}

We apply our principle for the derivation of modular equations to these
functions, and remark in addition the following:

The third roots of unity occurring in (16) are $=1$ if either $n$ is
indivisible by 3 and $\mu$ is divisible by 3, or $n''$ is divisible by 3
and $n'$ leaves the residue 2. In these cases every rational function of
$A$ and $B$ is a root of a transformation equation. In the other cases
this property belongs to the cube of such a rational function of $A$ and
$B$, for which the differences of the exponents of all the terms are
congruent to one another with respect to the modulus 3.

If these rational functions are \gesperrt{polynomials} and if
moreover all their values for $q = 0$ are finite, then they must be equal
to a constant, and thereby we obtain transformation equations.

If $n', n''$ are primes, then here too it suffices (cf. § 72) that the
negative powers of $q$ in the expansion of such a function in ascending
powers of $q$ drop out in the one principal case, namely in that where
\[
\begin{pmatrix} a, & 0 \\ c, & \partial \end{pmatrix}, \quad
\begin{pmatrix} a', & 0 \\ c', & \partial' \end{pmatrix}, \quad
\begin{pmatrix} a'', & 0 \\ c'', & \partial'' \end{pmatrix}
\]
are equal to
\[
\begin{pmatrix} 1, & 0 \\ 0, & n \end{pmatrix}, \quad
\begin{pmatrix} 1, & 0 \\ 0, & n' \end{pmatrix}, \quad
\begin{pmatrix} 1, & 0 \\ 0, & n'' \end{pmatrix},
\]
thus
\[
U = f(\omega)f(n'\omega)f(n''\omega)f(n'n''\omega);
\]
%CONT>>
